% notes-terms.tex
% Topic: Terms

\begin{definitionbox}{Definition ($\mathcal{L}$-Term)}
\begin{definition}[$\mathcal{L}$-Term]\label{def:l-term}
Let $\mathcal{L}$ be a first-order language with a fixed set of variables.
The \textbf{$\mathcal{L}$-terms} are defined recursively:
\begin{enumerate}
  \item every variable is an $\mathcal{L}$-term;
  \item every constant symbol of $\mathcal{L}$ is an $\mathcal{L}$-term;
  \item if $f$ is an $n$-ary function symbol of $\mathcal{L}$ and
        $t_1,\ldots,t_n$ are $\mathcal{L}$-terms, then
        $f(t_1,\ldots,t_n)$ is an $\mathcal{L}$-term.
\end{enumerate}
\end{definition}
\end{definitionbox}

\begin{remark*}[Standard quantified statement]
$\mathrm{Term}_{\mathcal{L}}$ is the least set containing all variables and
constant symbols of $\mathcal{L}$ and closed under function-symbol
application.
\end{remark*}

\begin{remark*}[Interpretation]
Terms are the formal expressions that name elements of a structure once the
variables have been assigned values.
\end{remark*}

\begin{dependencies}
\begin{itemize}
  \item \hyperref[def:first-order-language]{First-Order Language}
\end{itemize}
\end{dependencies}

\begin{definitionbox}{Definition (Closed Term)}
\begin{definition}[Closed Term]\label{def:closed-term}
A \textbf{closed term} is an $\mathcal{L}$-term containing no variables.
\end{definition}
\end{definitionbox}

\begin{remark*}[Standard quantified statement]
$t$ is closed iff $t\in\mathrm{Term}_{\mathcal{L}}$ and no variable occurs in
$t$.
\end{remark*}

\begin{remark*}[Interpretation]
Closed terms denote elements of a structure without requiring a variable
assignment.
\end{remark*}

\begin{dependencies}
\begin{itemize}
  \item \hyperref[def:l-term]{$\mathcal{L}$-Term}
\end{itemize}
\end{dependencies}

\begin{definitionbox}{Definition (Interpretation of Terms)}
\begin{definition}[Interpretation of Terms]\label{def:term-interpretation}
Let $\mathcal{M}$ be an $\mathcal{L}$-structure and let
$s\colon\mathrm{Var}\to M$ be a variable assignment. The
\textbf{value} $t^{\mathcal{M}}[s]$ of an $\mathcal{L}$-term $t$ in
$\mathcal{M}$ under $s$ is defined recursively by variables, constants, and
function-symbol application.
\end{definition}
\end{definitionbox}

\begin{remark*}[Standard quantified statement]
$x^{\mathcal{M}}[s]=s(x)$,
$c^{\mathcal{M}}[s]=c^{\mathcal{M}}$, and
\[
  f(t_1,\ldots,t_n)^{\mathcal{M}}[s]
  =
  f^{\mathcal{M}}(t_1^{\mathcal{M}}[s],\ldots,t_n^{\mathcal{M}}[s]).
\]
\end{remark*}

\begin{remark*}[Interpretation]
Term interpretation evaluates formal terms as actual elements of the universe
of a structure.
\end{remark*}

\begin{dependencies}
\begin{itemize}
  \item \hyperref[def:l-structure]{$\mathcal{L}$-Structure}
  \item \hyperref[def:l-term]{$\mathcal{L}$-Term}
\end{itemize}
\end{dependencies}

\begin{proposition}[Embeddings Preserve Term Interpretation]\label{prop:embeddings-preserve-term-interpretation}
Let $e\colon\mathcal{M}\to\mathcal{N}$ be an embedding of
$\mathcal{L}$-structures. For every $\mathcal{L}$-term $t$ and every variable
assignment $s\colon\mathrm{Var}\to M$,
\[
  e\bigl(t^{\mathcal{M}}[s]\bigr)
  =
  t^{\mathcal{N}}[e\circ s].
\]
\hyperref[prf:embeddings-preserve-term-interpretation]{\textit{Go to proof.}}
\end{proposition}

\begin{remark*}[Standard quantified statement]
$\forall e,\mathcal{M},\mathcal{N},t,s$, if $e$ is an embedding
$\mathcal{M}\to\mathcal{N}$, then
$e(t^{\mathcal{M}}[s])=t^{\mathcal{N}}[e\circ s]$.
\end{remark*}

\begin{remark*}[Predicate reading]
\[
\mathsf{Embedding}(e,\mathcal{M},\mathcal{N})
\Longrightarrow
e(t^{\mathcal{M}}[s])=t^{\mathcal{N}}[e\circ s].
\]
\end{remark*}

\begin{remark*}[Interpretation]
Embeddings commute with term evaluation: first evaluate in the smaller
structure and then embed, or first embed the assignment and then evaluate in
the larger structure.
\end{remark*}

\begin{dependencies}
\begin{itemize}
  \item \hyperref[def:l-structure-embedding]{Embedding of $\mathcal{L}$-Structures}
  \item \hyperref[def:term-interpretation]{Interpretation of Terms}
\end{itemize}
\end{dependencies}
